\addtocontents{toc}{\protect\setcounter{tocdepth}{2}}

\section{Applications}
\label{sec:applications}


Building on the established three-layer taxonomy (i.e. foundational, self-evolving and collective reasoning) mentioned in previous sections, we now examine how these capabilities manifest across real-world applications.
This section surveys representative reasoning-empowered agentic systems across several key domains, as illustrated in Figure~\ref{fig:application}, including math exploration and vibe coding (Section~\ref{sec:app-math}), scientific discovery (Section~\ref{sec:app-science}), robotics (Section~\ref{sec:app-embodied}),  healthcare (Section~\ref{sec:app-healthcare}), and autonomous web exploration and research (Section~\ref{sec:app-research}). Specifically, each domain exhibits distinctive forms of reasoning, influenced by its data modalities and environmental constraints.
Accordingly, our discussion in each subsection is organized around three layers:
(1) \textbf{core abilities} such as planning, tool use and search that span scientific hypothesis generation, embodied control, medical reasoning, automated experimentation and symbolic problem solving, for example;
(2) \textbf{self-evolving abilities} that integrate feedback, reflection and memory modules which refine domain-specific competence through iterative experiment loops, lifelong skill learning and clinical adaptation; and
(3) \textbf{collective multi-agent reasoning} that enables collaboration and specialization from cooperative scientific assistants to coordinated robotic teams, diagnostic ensembles or multi-aspect experts. This section highlights how agentic reasoning frameworks adapt to domain-specific knowledge structures and tasks, illustrating the transition from traditional LLM reasoning to goal-directed, domain-aware and active agentic intelligence. 



\subsection{Math Exploration \& Vibe Coding Agents}
\label{sec:app-math}
Mathematics and code have traditionally served as two of the most widely used domains for evaluating reasoning in artificial intelligence, as both require structured symbolic manipulation and precise multi-step deduction. Traditional benchmark-driven evaluation in these domains is showing clear limitations. Widely used math datasets such as GSM8K \cite{cobbe2021training}, MATH \cite{hendrycks2021measuring}, and AIME \cite{MAA_AIME} are increasingly saturated, which makes it difficult to distinguish among modern high-performing models. The problems in these datasets often rely on a small set of recurring techniques and do not require the sustained and exploratory reasoning needed to assess more advanced mathematical capabilities. Even recent evaluations such as FrontierMath \cite{glazer2024frontiermath} continue to emphasize final-answer accuracy, which offers only a partial view of an agent’s reasoning process and its ability to adjust strategies during problem solving.

Under the agentic reasoning paradigm, however, both areas are undergoing a substantial shift from static problem solving to dynamic processes that emphasize exploration, adaptation, and collaboration. In mathematics, recent systems \citep{novikov2025alphaevolve,luong2025towards,trinh2024solving,romera2024mathematical} demonstrate that agents can engage in competition-level reasoning, building on the success of LLMs in coding tasks. Work in foundational mathematics \citep{swirszcz2025advancing,georgiev2025mathematical} further shows that agents can search for new problems, propose conjectures, construct auxiliary lemmas, and explore deeper structures in mathematical concepts. These developments position mathematics not merely as an evaluation benchmark but as a domain of active \emph{mathematical exploration}.

Large Language Models have also reshaped coding through the emerging workflow known as agentic coding and \emph{vibe coding} \citep{karpathy2025vibecoding,willison2025vibecoding}. In this paradigm, the model acts as an interactive collaborator that engages in multi-turn natural-language dialogue. Users iteratively design and refine programs while the agent maintains context, adapts to evolving requirements, and continuously self-corrects. Modern tools such as Copilot\footnote{\url{https://github.com/features/copilot}} and Cursor\footnote{\url{https://cursor.com}} have further popularized this collaborative workflow, making interactive programming a common practice in real-world software development.

In this section, we organize our discussion according to the three-layer framework introduced earlier. The foundational layer (Section~\ref{sec:app-math-planning}) concerns the core reasoning and execution skills: mathematical agents perform symbolic manipulations and step-by-step derivations across arithmetic, algebra, geometry, and calculus, while code agents carry out syntax-aware generation, implement functions, and verify correctness through interpreter or compiler feedback. The self-evolving layer (Section~\ref{sec:app-math-evolving}) introduces mechanisms for reflection and adaptation. Mathematical agents learn from intermediate reasoning traces to correct missteps or explore alternative solution paths, and code agents iteratively debug, refine, and optimize implementations based on runtime feedback or test results. The collective layer (Section~\ref{sec:app-math-collaboration}) focuses on collaboration, where agents exchange intermediate results, share reusable modules, and jointly develop complex proofs or codebases. Taken together, these layers reveal how mathematics and coding are becoming domains in which agentic reasoning enables increasingly creative and adaptive problem solving.


\subsubsection{Foundational agentic reasoning}
\label{sec:app-math-planning}


\paragraph{Planning.}

Explicit planning is widely recognized as a core mechanism for enhancing the structured 
reasoning capabilities of LLMs. In the domain of mathematical discovery, several systems exhibit structures that can be interpreted as forms of planning. In representation theory and knot theory, the system of \citet{davies2021advancing} guides human mathematicians by proposing intermediate objects and promising avenues of exploration, which function as high-level suggestions for organizing problem-solving workflows. In geometric reasoning, \citet{trinh2024solving} solves Olympiad-level geometry problems by decomposing them into sequential stages of construction, lemma generation, and verification, yielding a structured multi-step process that resembles a planned reasoning trajectory. Program-search approaches \citep{romera2024mathematical} iteratively refine candidate programs and mathematical structures, a procedure that naturally forms a coarse-to-fine exploration path. Large-scale exploration frameworks \citep{georgiev2025mathematical,swirszcz2025advancing} also operate through cycles of proposing, testing, and modifying conjectures or geometric objects, which collectively create a procedural structure aligned with planning. Efforts toward more robust mathematical reasoning \citep{luong2025towards} similarly rely on stepwise reasoning patterns, further reinforcing the presence of implicit planning dynamics. of implicit planning dynamics across mathematical agents.

In code agents, planning has likewise emerged as an essential component for organizing 
multi-step reasoning and enabling more structured decision-making. Early systems such as 
CodeChain \citep{le2023codechain} and CodeAct \citep{wang2024executable} introduce explicit planning 
or action spaces to support modular code construction, while KareCoder \citep{huang2024knowledge} 
integrate external knowledge sources or domain-specific information 
into the planning process. Subsequent works explore more structured planning organizations, including 
multi-stage control flows \citep{bairi2024codeplan,multi-stage}, tree-shaped planning structures 
\citep{li2024codetree,ni2024tree}, and adaptive refinement mechanisms \citep{aggarwal2025dars}. 
Planning has also been linked to improved exploration breadth: GIF-MCTS \citep{dainese2024generating} 
incorporates Monte Carlo Tree Search to explore multiple code-generation trajectories.
Recent extensions demonstrate applicability in specialized domains such as hardware design, 
where VerilogCoder \citep{ho2025verilogcoder} employs graph-structured planning and waveform-based 
verification. To address environments where state serialization is difficult, Guided Search 
\citep{zainullina2025guidedsearchstrategiesnonserializable} introduces lookahead and trajectory 
selection strategies for evaluating candidate actions without full environment access.





\paragraph{Tool-Use.}
Integrating external computational tools with LLMs has become a central mechanism for extending the reasoning and generation capabilities of single-agent systems. A defining characteristic of many mathematical reasoning systems is their integration with external computational tools. Formal theorem-proving agents such as \citet{thakur2024context} operate directly within the Lean proof assistant, selecting tactics and interacting with the underlying prover through in-context guidance. Position papers on formal mathematical reasoning \citep{yang2024formal} emphasize that progress in mathematical AI will depend on systems that can call theorem provers, satisfiability solvers, and computer algebra systems as part of a broader reasoning loop. Program-search frameworks for discovery \citep{romera2024mathematical} rely on executing generated programs and employing symbolic routines for verification. Generative modelling approaches \citep{ellenberg2025generative} make use of computational number-theoretic tools to check and filter generated candidates. Geometry-focused systems \citep{trinh2024solving,hubert2024ai} integrate automated geometric solvers and checkers to validate constructions and derived relations. Across these systems, external computational resources play a central role in enabling correct and scalable mathematical reasoning.

In code agents, external tools have similarly become crucial for extending the capabilities of 
LLM-based agents beyond pure text generation. Early work such as Toolformer \citep{schick2023toolformer} 
and ToolCoder \citep{zhang2023toolcoder} explored how models can learn to invoke APIs or search tools 
to obtain missing information during generation. Subsequent systems integrate increasingly rich 
toolchains: ToolGen \citep{wang2024toolgen} leverages automatic completion tools to resolve undefined 
dependencies, while CodeAgent \citep{zhang2024codeagent} incorporates multiple programming utilities 
 including search, documentation reading, symbol navigation, and code execution to support more 
realistic software workflows. Several methods focus on improving tool-feedback loops, such as ROCODE 
\citep{jiang2024rocode}, which combines real-time error detection with adaptive backtracking, and 
CodeTool \citep{lu2025codetool}, which introduces process-level supervision to improve the reliability 
of tool invocation. Collectively, these systems show that tool integration provides essential external 
signals, via search results, documentation, static analysis, or execution feedback, that extend the 
reasoning and generation capabilities of single-agent LLMs.

\paragraph{Search and Retrieval.}

Search and retrieval has emerged as a complementary mechanism that enriches model contexts through external information sources.
Search is a recurring mechanism in mathematical discovery. Program-search based systems \citep{romera2024mathematical} treat mathematical discovery as navigating a program space in which candidate programs encode conjectures or structural hypotheses, with iterative filtering based on symbolic or numerical checks. Generative modelling approaches \citep{ellenberg2025generative} explore families of mathematical objects by sampling from flexible distributions that capture structural regularities. Geometric systems such as \citet{trinh2024solving} and \citet{swirszcz2025advancing} search over constructions, configurations, and high-dimensional polytopes, guided by learned heuristics or structural constraints. Large-scale discovery frameworks \citep{georgiev2025mathematical} operate through repeated propose--test--refine cycles across conjectures, supporting wide exploration over mathematical landscapes. All these systems rely on systematic search procedures that structure the exploration of mathematical ideas.

In code generation, repository-level retrieval systems 
such as RepoHyper \citep{phan2024repohyper} locate reusable code segments from large-scale code bases 
to provide more informative contexts for generation. CodeNav \citep{gupta2024codenav} dynamically 
indexes real repositories during generation, retrieving relevant functions and adjusting based on 
execution feedback. AUTOPATCH \citep{acharya2025optimizing} applies retrieval to performance 
optimization, combining historical code examples with control flow graph analysis for 
context-aware improvements. Structure-aware retrieval has also been explored: knowledge-graph-based 
repository representations \citep{athale2025knowledge} improve retrieval quality by capturing 
symbolic and relational structure, while cAST \citep{zhang2025cast} introduces AST-based chunking to 
enhance syntactic coherence and retrieval granularity. These retrieval methods demonstrate how 
external knowledge sources can augment single-agent LLMs by providing high-quality, structured 
contexts that guide both understanding and generation.

\subsubsection{Self-evolving agentic reasoning}
\label{sec:app-math-evolving}

\paragraph{Agentic Feedback and Reflection.}
Across mathematical and code reasoning tasks, feedback operates as an external signal that 
highlights discrepancies, confirms correct inferences, and directs the agent toward more reliable 
subsequent computations. Feedback mechanisms appear prominently across mathematical discovery systems. In program-search 
based discovery \citep{romera2024mathematical}, executing candidate programs and evaluating their 
outputs against constraints yields counterexamples or confirmations, enabling iterative refinement of 
conjectures. In geometry, automated checkers validate constructions and derived relationships 
\citep{trinh2024solving,hubert2024ai}, providing correctness signals that guide subsequent revisions. 
Interactive evaluation frameworks \citep{collins2024evaluating} show that human clarifications and 
follow-up prompts expose reasoning errors and improve model responses. Position work on formal 
reasoning \citep{yang2024formal} highlights verification, proof checking, and model checking as 
essential sources of structured feedback. In several systems involving multiple candidate hypotheses 
\citep{romera2024mathematical,ellenberg2025generative,georgiev2025mathematical}, 
the use of verification signals to retain promising candidates functions analogously 
to a fitness-based evaluation step, since these signals determine which hypotheses 
survive and which are discarded, thereby shaping the direction of subsequent exploration 
without introducing an explicit learning signal.

For code agents, feedback and reflection are central to improving reliability over multi-step 
reasoning. Fault-aware editing methods such as Self-Edit \citep{zhang2023self} incorporate 
execution-based signals to refine erroneous code, while Self-Repair 
\citep{olausson2024selfrepairsilverbulletcode} integrates code and feedback models to diagnose test 
failures and propose targeted corrections. More structured systems like LeDeX \citep{jiang2024ledex} 
combine stepwise annotation, execution-driven verification, and automated repair into a closed-loop 
pipeline in which feedback continually informs the next revision. Reflection also functions as a 
form of memory: iterative self-improvement frameworks such as Self-Refine \citep{madaan2023self}, 
Self-Iteration \citep{chang2023self}, and Self-Debug \citep{chen2023teaching} reuse earlier drafts, 
analyses, and explanations to guide subsequent revisions, while artifact-level mechanisms such as 
CodeChain \citep{le2023codechain} and LeDeX \citep{jiang2024ledex} retain reusable components, 
corrected snippets, and execution traces as persistent representations. Together, these approaches 
demonstrate how feedback—whether symbolic, execution-based, or self-generated—interacts with 
iterative memory to support structured refinement and long-horizon improvement in code-oriented 
agentic systems.


\paragraph{Memory.}
Memory provides agents with a mechanism for retaining and leveraging information from earlier 
reasoning steps, allowing them to maintain consistency, improve intermediate states, and improve 
their performance over extended problem-solving horizons. While few systems introduce an explicit memory module, many mathematical agents rely on forms of persistent state that can be viewed as implicit memory. Interactive evaluation frameworks \citep{collins2024evaluating} maintain conversational and problem-state context across multiple turns, allowing models to build upon earlier partial derivations. Formal-theorem-proving agents \citep{thakur2024context} operate over evolving proof states in Lean, which accumulate tactics, subgoals, and intermediate lemmas, functioning as structured persistent information. Program-search and discovery systems \citep{romera2024mathematical,georgiev2025mathematical} retain conjecture histories, counterexamples, and successful constructions as part of their iterative refinement processes. Their role in preserving and reusing information across reasoning steps aligns with the broader notion of memory in agentic systems.



In code agents, memory increasingly takes the form of explicit structures that maintain coherence
over long-horizon generation. Several systems construct shared or structured workspaces: 
Self-Collaboration \citep{dong2024self} introduces a blackboard memory for storing task descriptions,
intermediate drafts, and revision records, enabling agents to coordinate through a common 
representation. Architectural approaches such as L2MAC \citep{holt2023l2mac} and Cogito 
\citep{li2025cogito} extend this idea by organizing context into dedicated registers, hierarchical 
memory units, or long-term knowledge stores, overcoming context-window limits and supporting 
multi-file or large-function reasoning. Across these designs, the underlying insight is consistent: effective code agents require 
persistent, structured, and often domain-aware memory that preserves intermediate reasoning and 
enables self-improvement across extended development trajectories.





\subsubsection{Collective multi-agent reasoning}
\label{sec:app-math-collaboration}
To address the growing complexity of tasks in mathematical discovery and code generation, recent 
systems increasingly rely on multi-agent or modular designs that decompose problems into 
cooperating specialized components. Mathematical discovery frameworks often organize reasoning into explicitly defined 
multi-agent or multi-component workflows that collaborate to explore and validate mathematical ideas. The polytope-generation system \citep{swirszcz2025advancing} uses multiple specialized components that generate, evaluate, and refine geometric objects, forming a genuine collaborative workflow. Large-scale exploration frameworks \citep{georgiev2025mathematical} often divide discovery into modules for proposing conjectures, identifying counterexamples, and refining statements, which, although implemented within a unified system, mirror multi-agent role specialization. Early work on AI-assisted mathematical research \citep{davies2021advancing} and Olympiad-level systems \citep{hubert2024ai} also involve human--AI collaboration, where human mathematicians interact with AI systems in a complementary manner. These developments indicate that mathematical discovery is an inherently collaborative process, and multi-agent architectures provide a natural vehicle for expressing such collaboration in agentic systems.

Multi-agent systems for code generation have progressed from simple role-based pipelines to adaptive, 
collaborative frameworks capable of handling long-horizon software development. Early approaches such 
as Self-Collaboration \citep{dong2024self} and AgentCoder \citep{huang2023agentcoder} decompose tasks 
into sequential roles, while hierarchical designs like PairCoder \citep{zhang2024pair} and FlowGen 
\citep{lin2024soen} introduce an architecture in which high-level agents handle planning and lower-level 
agents carry out concrete implementation. Flexible systems such as SoA \citep{ishibashi2024self} further 
adjust the number and specialization of agents in response to task complexity. Other frameworks, 
including MapCoder \citep{islam2024mapcoder}, AutoSafeCoder 
\citep{nunez2024autosafecoder}, and QualityFlow 
\citep{hu2025qualityflow}, rely on repeated cycles in which multiple agents 
generate, test, analyze, and repair code. Recent work explores self-evolving system structures, as in 
SEW \citep{liu2025sew}, which reorganizes collaboration pathways based on runtime feedback, and EvoMAC 
\citep{hu2024self}, which adjusts agent strategies through an iterative text-based update mechanism. 
Collaborative optimization methods such as Lingma SWE-GPT \citep{ma2024lingmaswegptopendevelopmentprocesscentric}, 
CodeCoR \citep{pan2025codecor}, SyncMind \citep{guo2025syncmind}, and CANDOR 
\citep{xu2025hallucinationconsensusmultiagentllms} explicitly improve cross-agent coordination. 
Together, these systems show a clear shift toward multi-agent code generators that rely not only on 
role decomposition, but also on reflection, distributed evaluation, adaptive restructuring, and 
team-level optimization, transforming code generation into an increasingly coordinated and resilient 
problem-solving process.


\subsection{Scientific Discovery Agents}
\label{sec:app-science}
Scientific-discovery agents aim to accelerate the entire life cycle for scientific research, from hypothesis generation through experimental execution, by coupling LLMs with domain-specific simulators, laboratory automation and up-to-date literature. These systems ground decision in verifiable processes while handling heterogeneous data, safety constraints and long-horizon goals.  

In this subsection, we begin with the foundational layer (Section~\ref{sec:app-science-planning}), which encompasses planning under scientific context, tool-augmented interaction with scientific resources, search and retrieval mechanisms including RAG-based systems and execution-time integration with laboratory hardware.
Building upon these capabilities, the self-evolving layer (Section~\ref{sec:app-science-memory}) introduces agentic memory, feedback and reflection, which enable scientific agents to refine hypotheses, adapt protocols and learn from experimental outcomes.
Finally, the collective layer (Section~\ref{sec:app-science-multi}) explores multi-agent collaboration, where agents coordinate roles, share intermediate knowledge and jointly reason toward complex scientific goals.



\subsubsection{Foundational agentic reasoning}
\label{sec:app-science-foundational}


\paragraph{Planning.}
\label{sec:app-science-planning}


Scientific agents utilize reasoning-enhanced planning ability to decompose a research goal into steps, decides which tool or simulator to call next, then revises the plan as evidence arrives. In short, the \textit{chain of thought} emerges from LLM reasoning that compiles instructions into rigorous executable plans~\cite{wei2022chain}.
% There are certain ways that LLM reasoning interacts with agentic planning ability. 
For example, ProtAgents~\cite{ghafarollahi2024protagents} materializes a planner agent that utilizes LLM reasoning capability to formulate a concrete plan for protein analysis and keep modifying it with feedback from another critic agent, and Eunomia~\cite{ansari2023agentbasedlearningmaterialsdatasets} uses ReAct-style~\cite{yao2023react} workflow to make in-context reasoning: after retrieving a top-$k$ evidence set, the backbone LLM quote a warranting sentence, and that citation drives the next action choice. Other examples include MatExpert~\cite{ding2024matexpertdecomposingmaterialsdiscovery}, which deploys a chain-of-thought LLM to author a stepwise transition pathway and then emits a structured crystal candidate from a feedback loop.

Planning can also act as reasoning constraint. For instance, Curie~\cite{kon2025curierigorousautomatedscientific} utilizes a rigor engine to align, setup and do reproducibility check within planning steps proposed by the Architect LLM. Thus, the Architect's free-form reasoning cannot advance unless these rigor gates are satisfied, which transforms planning into both a guide and a regulator of the reasoning process. In addition, a general purpose biomedical agent, Biomni~\cite{huang2025biomni} constrains its reasoning within a dynamically constructed biomedical action space of comprehensive tools,  software packages and databases, requiring each hypothesis to be operationalized as executable code.


% ------------------------------------------------------------------ %
\paragraph{Tool-Use.}
\label{sec:app-science-tool}



Tool use is an important part of the reasoning loop for scientific agents nowadays. Specifically, rather than following rigid rules,   these agents can decide which tool and when to call, how to fill parameters and verify or revise based on evidence. For example, SciAgent~\cite{ma2024sciagenttoolaugmentedlanguagemodels} formalizes \textit{tool-augmented reasoning} as a four-step procedure: planning, retrieval, tool-based action and execution. Agents are trained to decide when to call a tool, which one, and how to integrate it into solving scientific tasks. Through domain-specific tools, ChemCrow~\cite{bran2023chemcrowaugmentinglargelanguagemodels} chains various expert chemistry tools so intermediate calculations become premises in the next reasoning step, which enables end-to-end planning and autonomous syntheses. CACTUS~\cite{mcnaughton2024cactus} similarly grounds explanations in cheminformatics outputs, reducing reliance on free-form reasoning by language models alone. 

Other notable examples include ChemToolAgent~\cite{yu2024chemtoolagent} and CheMatAgent~\cite{wu2025chematagentenhancingllmschemistry}. In particular, ChemToolAgent~\cite{yu2024chemtoolagent} employs a ReAct-like~\cite{yao2023react} architecture with multiple specialized chemistry tools, allowing the LLM to choose and parameterize tool calls while CheMatAgent~\cite{wu2025chematagentenhancingllmschemistry} pushes further by learning tool use: it integrates over 100 chemistry/materials tools, curates a tool-specific benchmark, and uses Monte Carlo Tree Search with step-level fine-tuning to learn both which tool to pick and how to fill arguments.

For biomedical agents, TxAgent~\cite{gao2025txagentaiagenttherapeutic} scales therapeutic reasoning across 211 vetted tools and it carries out multi-step reasoning that reconciles drug labels, interactions, and patient context—turning clinical justification into an executable trace. On the other hand, AgentMD~\cite{jin2024agentmdempoweringlanguageagents} builds a two-stage tool memory: it first mines thousands of clinical calculators from literature (i.e. making tools), then selects and applies the right ones at inference (i.e. using tools), pinning predictions to concrete computations. Other recent systems~\cite{lála2023paperqaretrievalaugmentedgenerativeagent, skarlinski2024languageagentsachievesuperhuman,chiang2024llamp,zhang2024honeycombflexiblellmbasedagent,qu2025crisprgptagenticautomationgeneediting,gao2025pharmagentsbuildingvirtualpharma} reinforce similar design: co-design tool-use with reasoning so each claim is computable and auditable.



Another notable category of tool-use is the ability of agentic execution, which includes but not limited to run codes and simulate environments. Execution layers bridge high-level plans to physical infrastructure, which enables scientific agents to autonomously operate laboratory hardware, orchestrate simulation pipelines, and manage large-scale data workflows. Recent works such as Organa~\cite{darvish2025organaroboticassistantautomated} ties LLM reasoning to task-and-motion planning plus scheduling and perception, executing multi-step experiments with autonomous robots;  AtomAgents~\cite{ghafarollahi2024atomagentsalloydesigndiscovery} exemplifies the simulation side of execution: a physics-aware system that plans and runs atomistic workflows, coordinating tools for code execution, analysis, and hypothesis checking; and Chemist-x~\cite{chen2025chemistxlargelanguagemodelempowered} shows wet-lab execution beyond a digital-only scenario, where agents generate control scripts and drive an automated platform to validate conditions without human intervention. 

Several other platforms couple execution with optimization or team-based autonomy. For instance, SGA~\cite{ma2024llmsimulationbileveloptimizers} formalizes the workflow of \textit{LLM-as-proposer and simulator-as-optimizer}  while MatExpert~\cite{ding2024matexpertdecomposingmaterialsdiscovery} operates a \textit{retrieval, transition and generation} workflow for material discovery tasks, and CellAgent~\cite{xiao2024cellagentllmdrivenmultiagentframework} coordinates planner, executor and evaluator roles to run full single-cell analysis pipelines. 



\paragraph{Search and retrieval.}

\label{sec:app-science-search}



Beyond simple context stuffing, recent scientific agentic systems elevate retrieval into a deliberate reasoning step: agents decide when and what to fetch, and how to use the evidence before committing to a hypothesis. With retrieval ability, BioDiscoveryAgent~\cite{roohani2025biodiscoveryagentaiagentdesigning} pulls literature and interim assay results inside a closed loop so the model’s next gene-perturbation choices are conditioned on what was read and measured; while DrugAgent~\cite{inoue2025drugagentexplainabledrugrepurposing} coordinates knowledge graph queries, targeted literature search through web API and machine learning predictors. Its planner selects retrieval actions and then reconciles heterogeneous evidence into an explainable rationale. To facilitate scientific research, ARIA~\cite{ramirezmedina2025acceleratingscientificresearchmultillm} operationalizes a \textit{search, filter then synthesis} workflow as role-bound steps that carry citations forward, turning literature into actionable procedures. Similarly, AI Scientist-v2~\cite{yamada2025ai} employs an agentic tree-search framework in which the agent actively queries scientific literature database during hypothesis formulation and manuscript drafting, ensuring that analyses and writing are grounded in existing evidence. For research idea generation, another recent work~\cite{qi2023largelanguagemodelszero} constrains the process with curated background packets, using retrieval as an experimental control.


Building on these developments, retrieval-augmented generation (RAG) frameworks position external sources not merely as supporting references but as active components of the reasoning process. Specifically, RAG-enhanced scientific agents make external sources as primary inputs to LLM context and reasoning material, mostly with explicit planning, passage extraction, citation and contradiction checks. For example, PaperQA~\cite{lála2023paperqaretrievalaugmentedgenerativeagent} and PaperQA2~\cite{skarlinski2024languageagentsachievesuperhuman} treat retrieval as the main loop. By deciding which documents to read, attributing every claim, and detecting conflicts to steer synthesis, these works can yield expert-level literature reviews that are inherently verifiable. In material science, LLaMP~\cite{chiang2024llamp} extends RAG beyond text. Specifically, it utilizes  hierarchical ReAct~\cite{yao2023react} agents call material-specific APIs to fetch band gaps or elastic tensors, edit structures and then reason with computed properties.



\subsubsection{Self-evolving agentic reasoning}
\label{sec:app-science-evolving}

Scientific discovery agents can go beyond static reasoning and acquire the ability to self-evolve, which is to learn from experience, refine their internal representations and improve decision quality over successive interactions. This self-evolving layer equips agents with mechanisms to monitor and revise their own reasoning, retain and reuse intermediate hypotheses and adjust future plans based on external feedback or environmental signals. In the following paragraphs, we discuss how memory modules enable the accumulation  of scientific knowledge and how feedback and reflection mechanisms support continual adaptation and reasoning consistency throughout long-horizon scientific workflows.


\paragraph{Memory.}
\label{sec:app-science-memory}

ChemAgent~\cite{tang2025chemagentselfupdatinglibrarylarge} implements a self-updating library. It decomposes chemistry problems into sub-tasks and writes reusable \textit{skills} (ex: procedures, patterns, solutions) that later prompts can retrieve and adapt, stabilizing long multi-step reasoning without re-deriving everything from scratch. On the other hand, MatAgent~\cite{takahara2025acceleratedinorganicmaterialsdesign} emphasizes interpretable generation for inorganic materials, where \textit{short-term memory} recalls recent compositions and feedback, \textit{long-term memory} preserves successful designs together with their reasoning traces, and both are reused across iterations to guide proposal refinement and enable transparent audit.

\paragraph{Agentic Feedback and Reflection.}
\label{sec:app-science-feedback-reflect}


Firstly, Scientific Generative Agent~\cite{ma2024llmsimulationbileveloptimizers} ties discrete LLM proposals to inner-loop simulations that optimize continuous parameters, advancing only when evidence improves. The reflection ability is driven by measurable loss reductions. Next, ChemReasoner~\cite{sprueill2024chemreasoner} performs heuristic search over the LLM’s idea space but scores and steers candidates with quantum-chemical feedback, turning electronic-structure signals into a principled critique of linguistic hypotheses. Complementing these physics-based signals, Curie~\cite{kon2025curierigorousautomatedscientific} embeds rigor check directly into control flow via intra-agent checks, inter-agent gates and an experiment-knowledge module. In parallel, LLMatDesign~\cite{jia2024llmatdesignautonomousmaterialsdiscovery} builds explicit self-reflection into materials workflows, prompting the agent to surface and repair inconsistencies before they propagate to tool calls. Moreover, NovelSeek~\cite{novelseekteam2025novelseekagentscientist} utilizes reflection as a closed loop, updating code and plans with human-interactive feedback after each round. Finally, a recent study~\cite{kumbhar2025hypothesisgenerationmaterialsdiscovery} regularizes the process up front with explicit goals \& constraints and afterwards with standardized scoring to provides an objective standard that makes reflection repeatable.



\subsubsection{Collective multi-agent reasoning}
\label{sec:app-science-multi}



Multi-agent frameworks for scientific discovery distribute labor across specialized LLM-driven roles, where advanced LLM reasoning not only orchestrates coordination between scientific agents but also adjudicates conflicting evidence to maintain coherence in the process. 

To illustrate, we introduce some important multi-agent frameworks as follows. Firstly, ProtAgents~\cite{ghafarollahi2024protagents} exemplifies this pattern in protein design. The framework involve agents for literature retrieval, structure analysis, physics simulation, and results analysis. Specifically, the backbone LLM directs reasoning over multi-modal outputs, choosing when to iterate or convergence-check based on feedback signals. PiFlow~\cite{pu2025piflow}, on the other hand, instantiates reasoning as principle-aware uncertainty reduction with a multi-agent loop in which a Planner agent relays strategy to a Hypothesis agent and a validation loop, explicitly tying multi-agent communication to hypothesis–evidence alignment.  AtomAgents~\cite{ghafarollahi2024atomagentsalloydesigndiscovery} also brings similar role specialization to alloy discovery. In particular, the agent uses LLM-guided reasoning to control over when to trigger simulations and how to evaluate multi-modal results, letting reasoning allocate computational resources and prune alloy candidates.



With a similar \textit{planner, executor and evaluator} framework, CellAgent~\cite{xiao2024cellagentllmdrivenmultiagentframework} instantiates researches on single-cell analysis, where the planner LLM reasoning selects tools or hyper-parameters and the evaluator LLM triggers self-iterative re-runs when quality checks fail. Some other notable works include ARIA~\cite{ramirezmedina2025acceleratingscientificresearchmultillm} that introduces a four-agent framework (scout, filter, synthesizer and procedure-drafter), Curie~\cite{kon2025curierigorousautomatedscientific} that embeds rigor into multi-agent planning, Team of AI-made Scientists (TAIS)~\cite{liu2024toward} for gene-expression discovery and the Virtual Lab~\cite{swanson2024virtual} for nano-body design with role agents. 

\subsection{Embodied Agents}
\label{sec:app-embodied}


Embodied agents extend reasoning beyond text, anchoring language in robotic perception, manipulation and navigation.  By embedding LLMs within robotic and simulated bodies, these embodied agents tackle real-world generalization, continual adaptation and multi-modal grounding.  

In this subsection, we begin with the foundational layer (Section~\ref{sec:app-embodied-planning}), which covers long-horizon embodied planning, tool-assisted perception, manipulation and execution. Building upon these capabilities, the self-evolving layer (Section~\ref{sec:app-embodied-evolve}) introduces agentic memory, feedback and self-reflection capabilities  enabling robots to refine control policies, adapt to novel environments and improve performance through continual interaction. Finally, the collective reasoning layer explores multi-robot collaboration (Section~\ref{sec:app-embodied-multi}), where agents coordinate perception, share learned representations and jointly reason about tasks to achieve complex embodied goals.

\subsubsection{Foundational agentic reasoning}


\label{sec:app-embodied-planning}

\paragraph{Planning.}


Early work such as SayCan~\cite{ahn2022icanisay} established the template by mapping linguistic descriptions to skill affordance estimates and  SayPlan~\cite{DBLP:conf/corl/RanaHGA0S23} refined this grounding by leveraging 3D scene graphs to align goal references with object-centric representations and spatial models. Beyond symbolic representations, EmbodiedGPT~\cite{mu2023embodiedgptvisionlanguagepretrainingembodied} use curated video CoT annotations of sub-goals to train models taht map multi-model input to structured sequences for embodied planning, while context-aware planning system~\cite{kim2024contextawareplanningenvironmentawarememory} adds semantic spatial map and object location information to the planning pipeline, enabling dynamical planning during execution.  In addition, DEPS~\cite{DBLP:journals/corr/abs-2302-01560} introduces an interactive planning loop (i.e. describe, explain, plan and select) for open-world multi-task agents. 


Embodied agents also rely on multi-modal reasoning traces that explicitly align perception with action. For example, Embodied CoT~\cite{zawalski2024robotic} trains vision-language-action models to generate reasoning steps incorporating visual features before executing an action. Fast\,ECoT~\cite{duan2025fast} accelerates this by caching and re-using reasoning segments across time-steps, reducing inference latency while preserving task success. More recently, Cosmos-Reason1~\cite{nvidia2025cosmosreason1physicalcommonsense} establishes an ontology of space, time and dynamics that lets CoT sequences encode structured physical priors. CoT-VLA~\cite{zhao2025cot} builds a visual chain-of-thought by predicting future image frames as intermediate sub-goals prior to action generation. Finally, Emma-X~\cite{sun2024emmaxembodiedmultimodalaction} integrates grounded chain-of-thought with look-ahead spatial reasoning, improving long-horizon embodied task performance.


Another line of works strengthen embodied planning through reinforcement learning, considering planning not only as static decomposition but as a self-evolving process that adapts to environment feedback. Robot-R1~\cite{kim2025robotr1reinforcementlearningenhanced} trains large VLMs to predict keypoint transitions under visual context, turning RL into a mechanism for learning physically grounded forward models. ManipLVM-R1~\cite{song2025maniplvm} exploits verifiable physical reward signals (e.g., trajectory match and affordance correctness) to reduce reliance on dense expert annotation. Embodied-R~\cite{zhao2025embodied} presents a collaborative framework where VLMs handle perception and smaller LMs handle reasoning, and the whole is trained via RL for embodied spatial reasoning. VIKI-R~\cite{kang2025viki} further extends this direction into heterogeneous multi-agent cooperation with a two-stage design, employing a two‐stage pipeline of chain-of-thought fine-tuning followed by hierarchical RL across agents coordinating activation and planning. 



\paragraph{Tool-use.}

\label{sec:app-embodied-tool}


Embodied agents can also be strengthened to interact with external tools to enhance perception and compensate for incomplete observations. GSCE~\cite{wang2025gscepromptframeworkenhanced}, for example, provides a prompt-framework that binds skill APIs and constraints for safe LLM-driven drone control. MineDojo~\cite{fan2022minedojo} links agents to internet-scale corpora and thus enabling richer affordance grounding. Physical AI Agents~\cite{bousetouane2025physicalaiagentsintegrating} further introduces a modular architecture and a retrieval augmented generation design pattern for embedding real-world physical interaction into LLM-driven agents. Beyond offline tool use, some systems treat the environment itself as an API. For example, Matcha agent~\cite{zhao2023chatenvironmentinteractivemultimodal} uses an LLM to issue queries about objects and scenes and thereby acquire perceptual information needed for task completion.



On the other hand, execution module is one of the most important tool type. It translates high-level language instructions into continuous motor commands, enabling embodied agents to act reliably in physical environments.  Early systems such as SayCan~\cite{ahn2022icanisay} uses language to invoke robot pick-and-place skills; while LEO~\cite{huang2024embodiedgeneralistagent3d} broaden execution to more general manipulation settings and Hi Robot~\cite{shi2025hi} uses a VLM reasoner to process complex prompts and a low-level action policy executes the chosen step. More recent efforts broaden the execution space: Gemini Robotics~\cite{geminiroboticsteam2025geminiroboticsbringingai} introduces a large-scale vision-language-action model for real-world robot control and Octopus~\cite{yang2024octopusembodiedvisionlanguageprogrammer} generates executable code in simulated environments that bridges planning and manipulation. 


Beyond single-agent control, hybrid pipelines couple reactive reflexes with language-guided policies to support complex domains. For example, CaPo~\cite{liu2024capo} incorporates an execution phase where agents carry out decomposed sub-tasks and adapt their meta-plan based on progress; COHERENT~\cite{liu2024coherent} embeds a robot executor module within its PEFA (i.e. proposal, execution, feedback and adjustment) loop, which ensures each assigned sub-task is acted and refined appropriately; and MP5~\cite{qin2024mp5multimodalopenendedembodied} integrates multi-modal perception to generate executable plans in open-ended Minecraft. At the perception–action interface, LLM-Planner~\cite{song2023llmplannerfewshotgroundedplanning} generates sub-goals and maps them into action sequences via a low-level controller and EmbodiedGPT~\cite{mu2023embodiedgptvisionlanguagepretrainingembodied} illustrate how LLM-generated plans can be translated into control policy for embodied control in physical environments. 



\paragraph{Search and retrieval.}

Embodied agents can also use search and retrieval ability to ground language in spatial structure and past experience. Early navigation systems such as L3MVN~\cite{yu2023l3mvn} use LLMs to query a semantic map and select promising frontiers as long-term goals during visual target navigation, while SayNav~\cite{rajvanshi2024saynav} and SayPlan~\cite{DBLP:conf/corl/RanaHGA0S23} build 3D scene graphs and then search task-relevant subgraphs so language instructions can be translated into grounded waypoints and sub-tasks in large environments. Long-horizon navigation works like ReMEmbR~\cite{anwar2025remembr} maintain a structured spatio-temporal memory that can be queried to answer “where” and “when” questions about past robot experience. Additionally, RAG-style systems make retrieval a first-class part of the planning loop: Embodied-RAG~\cite{quanting2024embodiedrag} and EmbodiedRAG~\cite{meghan2024embodiedrag} treat an agent’s experience and 3D scene graphs as non-parametric memories from which task-relevant episodes or subgraphs are retrieved for navigation and task planning; Retrieval-Augmented Embodied Agents~\cite{zhu2024retrieval} retrieve policies from a shared memory bank and condition action generation on them; and MLLM-as-Retriever~\cite{junpeng2024mllm} trains a multi-modal LLM retriever to rank past trajectories so each decision step can condition on the most useful prior experience rather than only the current observation.


\subsubsection{Self-evolving agentic reasoning}
\label{sec:app-embodied-evolve}



Embodied agents reliably achieve long-horizon autonomy when they can self-evolve over time: monitor their own internal states, store and update task-relevant knowledge and adjust behaviors when plans deviate. In the following paragraphs, we examine how memory modules, feedback signals and agentic reflection enable embodied agents to turn planning from a one-shot process into a continually improving cycle of behavior.


\paragraph{Memory.}
Effective memory mechanisms enable agents to reuse past experiences and maintain coherent task execution over extended interactions. Many systems cache recent observations in episodic buffers while summarizing long-term semantics in structured graphs, as in household planning~\cite{glocker2025llm} and long-horizon agents with hybrid multi-modal memory~\cite{li2024optimus1}. Skills and routines can be shared across tasks via indexed memory stores. For example, HELPER-X~\cite{sarch2023openendedinstructableembodiedagents} indexes discovered skills and action scripts, which aid future dialogue and can be shared across domains.
Spatial navigation methods such as BrainNav~\cite{ling2025endowingembodiedagentsspatial} maintain biologically inspired dual-map memories linked by a hippocampal hub to reduce hallucinations and drift. Broader contexts also benefit: CAPEAM~\cite{kim2024contextawareplanningenvironmentawarememory} incorporates environment-aware memory modules that track object states and spatial changes. Finally, lifelong episodic systems such as Ella~\cite{zhang2025ellaembodiedsocialagents} maintains long-term multi-modal memory system to support social-robot interaction.


\paragraph{Agentic Feedback and Reflection.}
Dialogue-based critique, calibrated uncertainty and environment-aware reward shaping refine policies beyond binary success signals. For example, Matcha agent~\cite{zhao2023chatenvironmentinteractivemultimodal} treats objects and scenes as interactive information sources before acting and FAMER~\cite{wang2025communicationefficientdesirealignmentembodied} uses lightweight preference feedback to adapt embodied agents to user intentions in real time. Uncertainty-aware planners such as KnowNo~\cite{ren2023robotsaskhelpuncertainty}, which proactively solicit guidance when confidence falls below guarantees, and Octopus~\cite{yang2024octopusembodiedvisionlanguageprogrammer}, which exploits environmental feedback to improve generated executable programs over time. At the multi-agent level, MindForge~\cite{licua2024mindforge} introduces theory-of-mind style perspective feedback so heterogeneous robots adapt to each other’s reasoning strategies; while ReAd~\cite{zhang2024efficientllmgroundingembodied} introduces a advantage-based feedback loop that enables an LLM planner to self-refine its collaboration strategies across embodied multi-agent tasks.



Robust reflection mechanisms help agents anticipate failures by monitoring their own reasoning and actions and then adjusting plans. Optimus-1~\cite{li2024optimus1} couples a \emph{Knowledge-guided Planner} with an \emph{Experience-Driven Reflector} to revise decisions using stored experience, while another recent study~\cite{Moncada_Ramirez_2025} defines structured agentic workflows (including self-Reflection, multi-Agent reflection and LLM Ensemble) that enable robots to reflect on and refine LLM-generated object-centered plans, thus reducing reasoning errors. Systems such as EMAC+~\cite{ao2025emac} interleave perception, planning and verification steps to perform online plan refinement and earlier works such as Voyager~\cite{wang2023voyager} also embeds an iterative prompting loop that uses environment feedback and execution errors to refine its skill library over time.


\subsubsection{Collective multi-agent reasoning}
\label{sec:app-embodied-multi}


Multi-agent collaboration enables embodied systems to divide labor and coordinate complex tasks more efficiently, with language often serving as the primary medium for negotiation and role allocation. For instance, SMART-LLM~\cite{kannan2024smartllmsmartmultiagentrobot} decomposes high-level instructions and allocates sub-tasks across multiple robots, while CaPo~\cite{liu2024capo} optimizes cooperative plans to avoid redundant exploration. For heterogeneity and coordination mechanisms, COHERENT~\cite{liu2024coherent} deploys a propose-execution-feedback-adjust loop across diverse robot types to enable seamless joint operation. In addition, Theory of Mind (ToM), which refers to an embodied agent’s ability to infer and reason about others' beliefs and mental states, is also highly related to embodied multi-agent systems~\cite{li2023theory, wu2025large, cross2024hypothetical}. For example, MindForge~\cite{licua2024mindforge} equips agents with explicit theory-of-mind representations and natural inter-agent communication to coordinate collaboratively.


For multi-modal frameworks, EMAC+~\cite{ao2025emac} integrate vision and language modules and continuously refine plans via visual feedback, COMBO~\cite{zhang2025combocompositionalworldmodels} integrates vision and language modules and continuously refine plans via visual feedback, and VIKI-R~\cite{kang2025viki} demonstrates reinforcement learning as a scalable coordination mechanism among embodied agents. At larger scales, studies such as RoCo~\cite{mandi2024roco} show how role negotiation and flexible protocols support adaptable teamwork in dynamic environments.

\subsection{Healthcare \& Medicine Agents}
\label{sec:app-healthcare}



Healthcare and medical agents seek to support the full clinical decision pipeline, from initial symptom triage to treatment planning and integrating LLMs with structured patient records, medical ontologies and expert guidelines. Unlike general assistants, these systems must operate under strict safety constraints, multi-modal evidence and legal justification.

In this subsection, we begin with the foundational layer (Section~\ref{sec:app-healthcare-foundational}), which includes medical and diagnostic reasoning and tool-augmented access to various biomedical knowledge bases and APIs. Building on these primitives, the self-evolving layer (Section~\ref{sec:app-healthcare-evolve}) examines memory, feedback and reflective modules that allow these agents to accumulate patient-specific context, adapt to longitudinal trajectories and revise clinical plans over time. Finally, the collective layer (Section~\ref{sec:app-healthcare-multi}) highlights multi-agent collaboration, which includes doctor–agent co-planning, human–AI shared autonomy and specialist model ensembles.



\subsubsection{Foundational agentic reasoning}
\label{sec:app-healthcare-foundational}



\paragraph{Planning.}


Planning is a core capability for healthcare agents, which enables them to structure long-horizon clinical pathways into diagnostic and treatment phases, refine workflows dynamically as patient conditions evolve and coordinate across teams and tools toward cohesive care delivery. We discuss several various recent advancements as follows. For instance, a recent agentic clinical system~\cite{ferber2024autonomousartificialintelligenceagents} orchestrates specialized tools and guideline citations to support oncology decision-making, EHRAgent \cite{wenqi2024ehragent} decomposes multi-table EHR inference into code-execution steps with feedback learning and PathFinder \cite{ghezloo2025pathfinder} presents a multi-agent, multi-modal histopathology workflow for diagnostic reasoning.  

Other frameworks model planning as an explicit orchestration layer across levels of abstraction. For example, MedAgent‑Pro \cite{wang2025medagent} proposes a hierarchical workflow which first generates disease-level diagnostic plans from guideline criteria and then dispatches tool-agent modules for execution. MedOrch \cite{he2025medorchmedicaldiagnosistoolaugmented} treats tool invocation itself as a planning primitive across modalities, orchestrating reasoning agents for multi-step diagnostic execution. On the other hand, ClinicalAgent \cite{yue2024clinicalagentclinicaltrialmultiagent} coordinates multi-agent workflows for clinical planning, leveraging LLM reasoning to allocate tools and synthesize evidence. In addition, planning in healthcare agents is increasingly adaptive, responding to new information and evolving contexts. For example, DoctorAgent‑RL \cite{feng2025doctoragentrlmultiagentcollaborativereinforcement} models clinical consultation as a dynamic decision-making process under uncertainty, optimizing questioning strategies and diagnostic paths via reinforcement learning; while DynamiCare \cite{shang2025dynamicaredynamicmultiagentframework} adjusts specialist-agent teams across multi-round interactions as new patient information emerges.



\paragraph{Tool-use.}


Tool integration significantly expands a healthcare agent’s action space, enabling precise calculations, medical image interpretation and access to specialized databases. Recent studies are summarized as follows. Several systems explicitly foreground extensibility. MedOrch~\cite{he2025medorchmedicaldiagnosistoolaugmented} introduces a modular architecture that allows new diagnostic APIs to be incorporated without retraining, while TxAgent~\cite{gao2025txagentaiagenttherapeutic} integrates over two hundreds pharmacological tools to support therapeutic decision-making across drug–disease–treatment relationships. AgentMD~\cite{jin2024agentmdempoweringlanguageagents} similarly curates and leverages over two thousands executable clinical calculators to learn risk-prediction pipelines.

Other approaches focus on structured function calling for safe execution. For example, LLM-based agents can reliably invoke bedside calculators when provided with explicit function signatures, ensuring arithmetic correctness in dosing and risk scoring~\cite{Goodell_2025}. MeNTi~\cite{zhu2025mentibridgingmedicalcalculator} goes further by enabling nested tool calls across multi-step medical calculators. Complementing these text-based integrations, MMedAgent~\cite{li2024mmedagent} demonstrates that agents can learn to select among multi-modal tools. 

In addition, execution is crucial for translating high-level clinical plans into concrete actions such as code operations, database queries or robotic procedures. VoxelPrompt~\cite{hoopes2024voxelpromptvisionlanguageagentgrounded} embed 3-D volumetric priors so that language instructions drive spatial segmentation and analysis of medical image volumes. On the other hand, embodied ultrasound-robot controllers~\cite{xu2024enhancingsurgicalrobotsembodied} translate LLM-generated plans into closed-loop robotic scanning via a “think-observe-execute” loop.
Adaptive reasoning-and-acting systems~\cite{dutta2024adaptivereasoningactingmedical} further refine both the reasoning and actions over time in simulated clinical environments. In medical imaging, systems like MedRAX~\cite{fallahpour2025medraxmedicalreasoningagent} materializes multi-step reasoning by integrating specialized chest-X-ray tools and LLM reasoning into an end-to-end diagnostic agent. PathFinder~\cite{ghezloo2025pathfinder} similarly executes multi-agent, multi-modal diagnostic workflows in histopathology.

Another class of healthcare agents deploys code-level workflows. For example, Conversational Health Agents~\cite{abbasian2024conversationalhealthagentspersonalized} compile dialogue actions into function calls and code execution for downstream processing, while EHRAgent~\cite{wenqi2024ehragent} materializes EHR operations via executable code. MedAgentGym~\cite{xu2025medagentgymscalableagentictraining} trains agents to produce code that is directly executed and graded, enforcing reliability of reasoning traces. DoctorAgent-RL~\cite{feng2025doctoragentrlmultiagentcollaborativereinforcement} validates multi-turn dialogue acts by executing reinforcement-learned strategies in simulated consultations, while AIPatient~\cite{yu2025simulated} materializes realistic patient scenarios for execution-based evaluation and another recent study~\cite{almansoori2025selfevolvingmultiagentsimulationsrealistic} demonstrates how self-evolving multi-agent simulations allow execution behaviors themselves to improve over time.



\paragraph{Search and retrieval.}
Search-based agents enhance clinical decision-making by linking LLM reasoning with external biomedical knowledge sources. For instance, MeNTi~\cite{zhu2025mentibridgingmedicalcalculator} supplements therapeutic reasoning by bridging LLM calls into multi-step medical calculators while EHRAgent~\cite{wenqi2024ehragent} dynamically executes code operations over multi-table EHR data to support complex tabular inference. Conversational Health Agents~\cite{abbasian2024conversationalhealthagentspersonalized} enrich personalized dialogue by integrating developer-defined external sources and orchestrating action flows. Another line of work explicitly embeds retrieval-augmented generation (RAG) into healthcare agents. For example, CLADD~\cite{lee2025ragenhancedcollaborativellmagents} retrieves molecular graphs and prior assay results before proposing compound hypotheses and MedReason~\cite{wu2025medreasonelicitingfactualmedical} issues targeted knowledge-graph sub-queries to anchor each reasoning step for clinical QA.




\subsubsection{Self-evolving agentic reasoning}
\label{sec:app-healthcare-evolve}



Self-evolving capabilities enable healthcare agents to maintain longitudinal clinical coherence. Representative use cases include accumulating relevant medical context across encounters, updating beliefs as new evidence arrives and revising decisions when inconsistencies surface. In the following paragraphs, we examine how memory, feedback and reflective mechanisms collectively turn clinical reasoning from a one-shot prediction into a continually improving care process.


\paragraph{Memory.}
Persistent memory is essential for tracking medical or patient history and maintaining context across interactions. For instance, epidemic-modeling agents~\cite{williams2023epidemicmodelinggenerativeagents} maintain temporal contact histories to trace infection chains over time; while MedAgentSim~\cite{almansoori2025selfevolvingmultiagentsimulationsrealistic} stores experience histories and refine diagnostic strategies over time. In structured data settings, EHRAgent~\cite{wenqi2024ehragent} records intermediate computations over tabular EHRs so subsequent steps can reference prior results. EvoPatient~\cite{du2025llmssimulatestandardizedpatients} interleaves memory with coevolution maintains evolving clinical state across dialogue phases while AIPatient~\cite{yu2025simulated} persists longitudinal EHR-derived variables to drive consistent responses. Multi-agent systems such as MedOrch~\cite{he2025medorchmedicaldiagnosistoolaugmented} contain clinical knowledge graph agent which can be considered as external memory that can be queried to retrieve known relationships or diagnostic patterns. 



\paragraph{Agentic Feedback and Reflection.}
Agentic feedback and self-reflection are complementary mechanisms that improve reliability and adaptability of healthcare agents. Feedback converts execution outcomes into learning signals: DynamiCare~\cite{shang2025dynamicaredynamicmultiagentframework} updates multi-agent treatment strategies when newly observed patient state contradicts prior plans; DoctorAgent-RL~\cite{feng2025doctoragentrlmultiagentcollaborativereinforcement} optimizes questioning policies from consultation rewards; and MedAgentGym~\cite{xu2025medagentgymscalableagentictraining} enforces correctness by executing and grading generated code. Tool-use pipelines also propagate execution feedback. For example, the success/failure of table queries in EHRAgent~\cite{wenqi2024ehragent} or calculator calls in MeNTi~\cite{zhu2025mentibridgingmedicalcalculator} and clinical-calculation agents~\cite{Goodell_2025} to refine subsequent actions.





\subsubsection{Collective multi-agent reasoning}
\label{sec:app-healthcare-multi}


Multi-agent collaboration is central to healthcare AI, since clinical decision-making often depends on consensus among specialists, negotiation of competing hypotheses and coordination across roles such as physicians, patients and trial designers. In the following, we discuss several strands of research centered around multi-agent capabilities.



For collaborative decision-making, notable frameworks include MDAgents~\cite{kim2024mdagents}, which automatically assigns tailored collaboration structures to teams of LLMs depending on medical task complexity, and DoctorAgent-RL~\cite{feng2025doctoragentrlmultiagentcollaborativereinforcement}, which uses a multi-agent reinforcement-learning framework to optimize multi-turn doctor-patient consultation dialogues. In addition, Agent-derived Multi-Specialist Consultation (AMSC)~\cite{wang2024directdiagnosisllmbasedmultispecialist} explores staged multi-specialist dialogues for differential diagnosis that mimics the medical scene of a patient consulting with multiple specialists. Other notable works include ClinicalAgent~\cite{yue2024clinicalagentclinicaltrialmultiagent}, which organizes clinical trial workflows via role-based agent collaboration / LLM reasoning and PathFinder~\cite{ghezloo2025pathfinder}, which integrates a diverse set of agents that can gather evidence and provide comprehensive diagnoses with natural language explanations. 

On the other hand, there are studies focusing on simulation-driven collaboration. These works highlight how multi-agent setups enrich training and evaluation. MedAgentSim~\cite{almansoori2025selfevolvingmultiagentsimulationsrealistic} co-evolves doctor and patient agents to simulate real-world multi-turn clinical interactions, and EvoPatient~\cite{du2025llmssimulatestandardizedpatients} uses co-evolution of patient and doctor agents to generate diagnostic dialogue data and therefore gathers experience to improve the quality of both questions and answers to enable accurate human doctor training. In addition, DynamiCare~\cite{shang2025dynamicaredynamicmultiagentframework} initiates a team of specialist agents that iteratively queries the patient system to integrate new information and adapts the composition and strategy. Finally, medical agents can also collaborative to aid medical reasoning process. For example, MedAgents~\cite{xiangru2023medagents} demonstrates zero-shot cooperation among domain-specialist agents in medical reasoning tasks, CLADD~\cite{lee2025ragenhancedcollaborativellmagents} uses retrieval-augmented generation to support drug-discovery workflows across agents and GMAI-VL-R1~\cite{su2025gmaivlr1harnessingreinforcementlearning} combines multi-modal reasoning and reinforcement learning in a multi-agent framework to support large-scale medical decision-making.



\subsection{Autonomous Web Exploration \& Research Agents}
\label{sec:app-research}

Web agents, GUI agents and autonomous research agents constitute three interlinked but distinct trajectories of agentic reasoning systems. Firstly, web agents specialize in navigating online resources, issuing web API calls or browser actions to retrieve dynamic evidence and steer research direction. GUI agents go further by manipulating software interfaces and multi-modal dashboards directly (i.e. clicking, typing, navigating) to execute experiments, data workflows and interface-based tasks. Autonomous research agents sit at the top of this hierarchy, pairing LLM reasoners with scientific workflows, tool-chains and meta-loops to drive hypothesis generation, data synthesis and paper writing. The core connection is a progression of autonomy: first web agents retrieve evidence from online resources, then GUI agents operationalize actions inside software interfaces, and finally autonomous research agents orchestrate full scientific workflows end-to-end.

In this subsection we begin with the foundational layer (Section \ref{sec:app-auto-foundational}), which captures the core capabilities that any autonomous agent must support: perceiving its environment, reasoning about goals, planning actions and grounding those into tool-augmented workflows. Building on these primitives, the self-evolving layer (Section \ref{sec:app-auto-evolving}) examines how agents incorporate feedback, memory and reflection to iteratively refine their behaviors and improve methods over time. Finally, the collective layer (Section \ref{sec:app-auto-multi-agent}) highlights how agents move beyond individual competence into coordination, specialization and emergent collaboration. While web agents, GUI agents and autonomous agents share common themes of goal-directed autonomy, tool-use and iterative improvement, they differ in where they act on, how they manipulate their environment and what goal they aim to achieve.



\subsubsection{Foundational agentic reasoning}


\label{sec:app-auto-foundational}


\paragraph{Planning.}



Planning is essential for web agents because they must decompose long-horizon tasks into manageable steps, adapt to dynamic pages and coordinate tool/invocation strategies. Early work such as WebGPT~\cite{nakano2021webgpt} fine-tuned GPT-3~\cite{brown2020language} to answer open-ended questions via a text-based web-browser interface. Then, various web-based methods deepened the planning paradigm: for example, SEEACT~\cite{zheng2024gpt} explored large multi-modal models as generalists that integrating visual and HTML grounding for web-based tasks, and AutoWebGLM~\cite{lai2024autowebglm} introduced HTML simplification and various learning techniques for open-domain web task decomposition and navigation. These works paved the way for recent systems such as Agent Q~\cite{putta2024agent} that integrate guided MCTS, self-critique and off-policy preference optimization on web-task benchmarks, and set the stage for even more advanced long-horizon web planners such as WebExplorer~\cite{liu2025webexplorer} and WebSailor~\cite{li2025websailor}.




% rl for web agents
In addition, reinforcement learning has become a core tool for improving the decision-making and planning behavior of web-based LLM agents. WebRL~\cite{qi2024webrl} introduces a self-evolving online curriculum that generates new tasks from unsuccessful attempts and trains an outcome-supervised reward model to guide policy optimization. WebAgent-R1~\cite{wei2025webagent} performs end-to-end multi-turn RL, learning web interaction policies directly from online rollouts with binary success rewards. DeepResearcher~\cite{zheng2025deepresearcher} scales RL to real-world web environments, using a multi-agent browsing architecture and exhibiting emergent behaviors such as plan formulation, cross-source corroboration, and self-reflection. Hybrid pipelines like AutoWebGLM~\cite{lai2024autowebglm} combine supervised training with RL fine-tuning to strengthen task decomposition and structured navigation, while Navigating WebAI~\cite{thil2024navigating} combines supervised learning and RL techniques to improve web navigation performance. Methods such as Pangu DeepDiver~\cite{shi2025pangu}, EvolveSearch~\cite{zhang2025evolvesearch}  and WebEvolver~\cite{fang2025webevolver} use RL-based self-improvement, for example by adaptively scaling search depth or jointly training an agent and a world-model-like simulator to improve long-horizon web decision-making. Hierarchical approaches like ArCHer~\cite{zhou2024archer} optimize high-level and low-level policies with a multi-turn hierarchical RL framework, while PAE~\cite{zhou2025proposer} combines a task proposer, an acting agent and an evaluator to support autonomous skill discovery via RL in internet environments. 


% gui agents



% normal planning for gui agents
Planning is a core capability for GUI agents, enabling them to coordinate long, multi-step interactions across applications and operating environments. OS-Copilot~\cite{wu2024copilot} approaches this by treating the desktop as a unified control space in which a generalist agent continually refines its multi-step workflows. Agent S~\cite{agashe2024agent} builds an experience-augmented planning stack that decomposes tasks into sub-goals while retrieving past trajectories and external knowledge to guide action sequencing. InfiGUIAgent~\cite{liu2025infiguiagent} strengthens planning by integrating hierarchical task structuring into a multi-modal backbone, allowing agents to organize GUI procedures at multiple levels of abstraction. MobA~\cite{zhu2025moba} and PC Agent~\cite{liu2025pc} employ hierarchical architectures that separate high-level planning from low-level execution—on mobile and desktop respectively. GUI foundation models such as OS-ATLAS~\cite{wu2024atlas}, OSCAR~\cite{wang2024oscar} and UItron~\cite{zeng2025uitron} further emphasize robust cross-application planning: OS-ATLAS offers a platform-agnostic action model for consistent control, OSCAR maintains state-aware plans that adapt as execution unfolds and UItron unifies offline and online planning within a single general-purpose GUI agent. 


Likewise, reinforcement learning has become a central way to endow GUI agents with planning over long action sequences. End-to-end frameworks such as ARPO~\cite{fanbin2025arpoendtoend} and ComputerRL~\cite{hanyu2025computerrl} directly optimize multi-step GUI trajectories with replay buffers or large-scale online interaction, replacing hand-crafted scripts with learned policies for general desktop control. R1-style and semi-online methods, including UI-R1~\cite{zhengxi2025uir1}, GUI-R1~\cite{run2025guir1}, InfiGUI-R1~\cite{yuhang2025infiguir1} and UI-S1~\cite{zhengxi2025uis1}, start from strong vision–language backbones and then use RL to sharpen action prediction and long-horizon reasoning. A complementary line focuses on where to act by improving visual grounding: GUI-Bee~\cite{yue2025guibee}, SE-GUI~\cite{xinbin2025enhancing}, UIShift/GUI-Shift~\cite{gao2025guishiftenhancingvlmbasedgui} and UI-AGILE~\cite{shuquan2025uiagile} develop RL-based grounding frameworks to help agents reliably localize target elements before executing actions. ZeroGUI~\cite{chenyu2025zerogui} pushes toward fully automated online RL loops, where the agent generates its own tasks and trajectories and improves with zero human annotation, while ComputerRL~\cite{hanyu2025computerrl} scales end-to-end online RL in large distributed desktop environments. AgentCPM-GUI~\cite{zhong2025agentcpmgui} couples supervised pre-training with reinforcement fine-tuning to strengthen decision quality on mobile apps. Finally, foundation-style GUI agents such as AutoGLM~\cite{xiao2024autoglm} and Mobile-Agent-v3~\cite{jiabo2025mobileagentv3} serve as general backbones that unify perception, grounding and action, and are trained or fine-tuned with scalable RL frameworks to align long-horizon GUI planning with real-world success signals.

% research agents
For autonomous research agents, planning modules translate abstract goals into actionable research itineraries.  For example, Agent Laboratory~\cite{schmidgall2025agentlaboratoryusingllm} organizes work into three structured stages, namely literature review, experimentation and report writing, and supports the workflow with tool-hooks that automate code execution, experiment runs and documentation. GPT Researcher~\cite{Elovic2023gptresearcher} uses a plan → research → write cycle, where a dedicated planner drafts the outline, retrieval/analysis agents gather evidence and a writer compiles the final report. Chain of Ideas~\cite{li2024chainideasrevolutionizingresearch} retrieves literature into a chain structure to reflect domain progression and support ideation via experiment design, whereas IRIS~\cite{garikaparthi2025iris} performs hypothesis exploration via Monte Carlo Tree Search to expand promising branches before committing to downstream tasks. Broader variants include ARIA~\cite{ramirezmedina2025acceleratingscientificresearchmultillm} and NovelSeek~\cite{novelseekteam2025novelseekagentscientist} that automate the research workflow with a complete literature search, hypothesis generation and experiment planning cycle.




\paragraph{Tool-Use.}



% web agents
For web agents, tool-use abilities underpins execute plans in realistic, dynamic environments. For example, WebVoyager~\cite{he2024webvoyager} systematizes multi-modal execution by building an end-to-end agent that operates on real websites. On the interaction side, BrowserAgent~\cite{zhang2025browseragent} makes the action space more human-like, defining a compact set of browser primitives (e.g., click, scroll, type) and coupling them with an explicit memory mechanism to maintain key conclusions across steps, yielding strong gains on multi-hop QA benchmarks. Finally, methods like WALT~\cite{prabhu2025walt} and pipeline-oriented systems such as WebDancer~\cite{wu2025webdancer} and WebShaper~\cite{tao2025webshaper} push tool use from mere execution toward tool discovery and data-centric interaction. Specifically, WALT teaches agents to reverse-engineer reusable tools from website functionality, while WebDancer and WebShaper embed web actions inside multi-turn information-seeking and dataset-synthesis loops, respectively.


% gui agent tool-use
Tool use is another core capability for GUI agents, enabling them to invoke system functions and application features as structured tools. As pioneering systems, AutoDroid~\cite{wen2024autodroid} automatically analyzes Android apps to construct functionality-aware UI abstractions that LLM agents can reason over as capabilities rather than raw layouts, while its successor AutoDroid-V2~\cite{wen2025autodroid} re-frames mobile UI automation as LLM-driven code generation, with an on-device small language model emitting executable scripts for a local interpreter. MobileExperts~\cite{zhang2024mobileexperts} models each expert as a tool-capable specialist and uses a dual-layer controller to select which expert and its associated tool-set to invoke at different stages of a mobile workflow. AgentStore~\cite{jia2025agentstore} pushes this idea to the platform level by treating heterogeneous agents themselves as tools: a MetaAgent uses AgentTokens to route operating-system subtasks to the most suitable specialized “tool-agent” through a unified interface. OS-Copilot~\cite{wu2024copilot} and OSCAR~\cite{wang2024oscar} integrate rich system-level tools into unified computer-control frameworks, so that complex desktop tasks are expressed as sequences of tool calls. OS-ATLAS~\cite{wu2024atlas} complements these systems with a foundation action model that offers robust cross-platform GUI grounding, serving as a reliable actuator layer for downstream tool-using agents. Finally, SeeClick~\cite{cheng2024seeclick} strengthens the execution stack by pre-training a visual GUI agent for GUI grounding, improving the ability to locate the correct on-screen elements from instructions.


Specialized tools can expand an autonomous research agent’s capabilities beyond pure text, allowing more fine-grained ability. For instance, Agentic Reasoning~\cite{wu2025agentic} automatically routes queries to appropriate tool modules like code execution, web search and structured memory agents when the main LLM detects a gap in reasoning; while Webthinker~\cite{li2025webthinkerempoweringlargereasoning} empowers autonomous web exploration and page navigation during long-horizon investigations, by interleaving reasoning, search and draft-writing with a web explorer module. PaperQA~\cite{lála2023paperqaretrievalaugmentedgenerativeagent} and its follow-up synthesis agent~\cite{skarlinski2024languageagentsachievesuperhuman} integrate PDF parsing and citation-level grounding to produce verifiable answers and literature syntheses, while Scideator~\cite{radensky2025scideatorhumanllmscientificidea} provides an IDE-style tool-chain that combines paper facets with novelty checks for real-time brainstorming. In addition, DeepResearcher~\cite{zheng2025deepresearcher} shows that reinforcement learning over real-web interactions improves deep-research efficiency and quality, with emergent behaviors such as plan refinement and cross-source corroboration. 


Execution components ground high-level reasoning in code, simulations or laboratory protocols to produce verifiable scientific outcomes. Agent Laboratory~\cite{schmidgall2025agentlaboratoryusingllm} executes experiments specified in declarative configuration files by orchestrating external toolchains, while Agentic Reasoning~\cite{wu2025agentic} integrates a coding agent that executes Python alongside web search and structured memory, feeding the results back into the reasoning process. MLR-Copilot~\cite{li2024mlrcopilotautonomousmachinelearning} turns research plans into runnable implementations via an ExperimentAgent that leverages retrieved prototype code, runs experiments, and iteratively debugs implementations. Dolphin~\cite{yuan2025dolphin} closes the loop by generating ideas, implementing them through code templates with traceback-guided debugging, executing experiments, and using the analyzed results to steer the next research cycle. The AI Scientist~\cite{chris2024the} automates end-to-end ML experiments, i.e. generating ideas, writing code, executing experiments and visualizing results, so that observed outcomes can guide subsequent runs, while The AI Scientist-v2~\cite{yamada2025ai} adds a dedicated experiment manager and progressive agentic tree search to prioritize and schedule experiment branches. Most recently, NovelSeek~\cite{novelseekteam2025novelseekagentscientist} introduces a unified closed-loop multi-agent framework that spans hypothesis generation, idea-to-methodology construction, and multi-round automated experiment execution with feedback across diverse scientific domains.


\paragraph{Search and Retrieval.}


Search and retrieval lie at the heart of what differentiates web agents from static language models: they must locate, synthesize and refactor web-scale information in dynamic environments. WebExplorer~\cite{liu2025webexplorer} tackles this by generating challenging information-seeking trajectories and training agents to interleave a search tool and a browse tool over many turns, resulting in improved multi-step retrieval policies on complex benchmarks. WebSailor~\cite{li2025websailor} likewise focuses on information-seeking under extreme uncertainty, constructing high-uncertainty search tasks and using a two-stage post-training pipeline  to instill uncertainty-reducing search strategies for long-horizon web tasks. INFOGENT~\cite{reddy2025infogent} also performs  multi-query search across diverse web sources, enabling comprehensive information retrieval beyond task completion. For retrieval-augmented generation applications, RaDA~\cite{kim2024rada} explicitly disentangles web-agent planning into Retrieval-augmented Task Decomposition and Retrieval-augmented Action Generation, so that each high-level subgoal and concrete action is conditioned on fresh search results while respecting context limits. In addition, GeAR~\cite{shen2025gear} advances retrieval itself by augmenting a base retriever with graph expansion and an agent framework, enabling multi-hop passage retrieval along graph-structured evidence chains. Finally, WebRAGent~\cite{anonymous2025webragent} exemplifies retrieval-augmented generation for web agents by retrieving past trajectories and external knowledge into a multi-modal RAG policy.



% gui agents
Several GUI agents use retrieval capability to inject external experience or knowledge at inference time. Synapse~\cite{zheng2023synapse} maintains an exemplar memory of abstracted trajectories and, for each new task, retrieves similar past trajectories as in-context plans, substantially improving multi-step decision-making. LearnAct~\cite{liu2025learnact} builds a three-agent pipeline that mines human demonstrations into a knowledge store and retrieves the most relevant instructions to guide mobile GUI execution on unseen and diverse tasks. MobileGPT’s Explore–Select–Derive–Recall framework~\cite{lee2023explore} equips a phone agent with human-like app memory, storing modular procedures that can be recalled and recomposed when similar tasks reappear. TongUI~\cite{zhang2025tonguiinternetscaletrajectoriesmultimodal} turns large-scale multi-modal web tutorials into the GUI-Net trajectory corpus, effectively giving agents a large offline memory of how humans operate hundreds of apps across multiple operating systems. RAG-GUI~\cite{xu2025retrieval} makes retrieval explicit at inference time by querying web tutorials and generating textual guidelines that are fed into any VLM-based GUI agent as step-by-step hints. WebRAGent~\cite{anonymous2025webragent} shows a related pattern in web automation, combining a multi-modal retriever with a web agent so that each action is conditioned on retrieved guidance.


% research agents
Search modules probe the research landscape to surface relevant papers and passages, enriching context and grounding subsequent reasoning. WebThinker~\cite{li2025webthinkerempoweringlargereasoning} equips large reasoning models with a Deep Web Explorer module for autonomous web search and page navigation, and uses an Autonomous Think–Search–and–Draft strategy with RL-based training to decide when to browse and what to extract during long-horizon tasks. DeepResearcher~\cite{zheng2025deepresearcher} scales end-to-end training on the real web via reinforcement learning in live search environments, optimizing the iterative think–search loop and exhibiting emergent behaviors such as plan formulation, cross-source corroboration, and self-correction over multi-step research trajectories. Retrieval-centric agents like PaperQA~\cite{lála2023paperqaretrievalaugmentedgenerativeagent} and its successor PaperQA2~\cite{skarlinski2024languageagentsachievesuperhuman} demonstrate that tightly coupling full-text retrieval with generation can substantially improve scientific QA accuracy while preserving cited provenance for literature synthesis.

In research settings, retrieval-augmented generation (RAG) grounds idea generation and analysis in freshly retrieved and citable passages. For example, GPT Researcher~\cite{Elovic2023gptresearcher} is an autonomous research-agent that retrieves sources and generates reports with citations, enabling traceability of claims to evidence. Chain of Ideas~\cite{li2024chainideasrevolutionizingresearch} organizes relevant literature into a chain-structured scaffold that mirrors a field’s progressive development, thereby guiding retrieval and ideation toward subsequent links in the argument. Meanwhile, Scideator~\cite{radensky2025scideatorhumanllmscientificidea} extracts key paper facets (e.g., purpose, mechanism, evaluation) and leverages them to drive targeted retrieval and recombination of ideas for identifying methodological or evidentiary gaps.




\subsubsection{Self-evolving agentic reasoning}
\label{sec:app-auto-evolving}



Effective self-evolving abilities enable these autonomous agents to adapt their behavior over time, retain crucial task context across interaction cycles and incrementally refine planning and execution strategies. The following paragraphs review how memory, feedback and self-reflection mechanisms support this continual improvement across these agent families, turning interaction from a one-shot pipeline into an iterative learning loop.





% web agents
\paragraph{Memory.}


Memory modules transform brittle, single-pass web interactions into reusable experience. For example, Agent Workflow Memory (AWM)~\cite{Wang2024AgentWM} induces reusable workflows from successful trajectories and retrieves them to guide future tasks, while ICAL~\cite{sarch2024vlm} distills noisy trajectories into high-level verbal and visual abstractions that are stored as a memory of multimodal experience and later injected into prompts. Control-oriented designs such as BrowserAgent~\cite{zhang2025browseragent} maintain explicit histories of past actions and intermediate conclusions in the agent’s context, instead of only re-encoding the current page view. GLM-based agents like AutoWebGLM~\cite{lai2024autowebglm} and AgentOccam~\cite{yang2024agentoccam} emphasize compressed page representations, using HTML simplification and carefully tuned observation spaces so that the agent’s prompt contains a shorter, more informative view of the state, with past steps preserved through the usual action–observation history. More integrated frameworks like LiteWebAgent~\cite{zhang2025litewebagent} expose planning, memory and tree search as modular components, and can plug in workflow memories together with search traces for long-horizon reuse.



% gui agents
Recent GUI agents adopt explicit memory modules that store and retrieve task-relevant information during long-horizon execution. Earlier work such as MobileGPT~\cite{lee2024mobilegpt} equips a mobile assistant with human-like app memory: it decomposes procedures into modular sub-tasks that are explored, selected, derived, and then stored so they can be recalled and reused instead of being re-discovered from scratch. Chain-of-Memory (CoM)~\cite{gao2025chain} incorporates short- and long-term memory by recording action descriptions and task-relevant screen information in a dedicated memory module, enabling cross-application navigation to track task state. More recent systems build increasingly structured memories: MobA’s multifaceted memory module~\cite{zhu2025moba} maintains environment- and user-level traces that an adaptive planner retrieves when refining mobile task plans, while MGA~\cite{cheng2025mga} represents each step as a triad of current screenshot, spatial layout, and a dynamically updated structured memory that summarizes past transitions, mitigating error accumulation in long chains of actions. Mobile-Agent-E~\cite{wang2025mobileagente} adds a persistent long-term store of tips and shortcuts distilled from prior trajectories, so later plans can call reusable guidance and subroutines instead of relearning them. Mirage-1~\cite{xie2025mirage} similarly organizes experience into a hierarchical skill memory that a planner can retrieve as reusable building blocks for new GUI tasks. 



% research agents
Long-term memory is crucial for autonomous research agents because it enables accumulation and reuse of prior knowledge, fostering continuity across research cycles. For example, Agent Laboratory~\cite{schmidgall2025agentlaboratoryusingllm} retains prior experiment code, results, and interpretation across its multi-phase workflow, enabling later stages to build on earlier work. GPT Researcher~\cite{Elovic2023gptresearcher} generates reports with embedded citations and provides context for planning and extension of research topics. Chain of Ideas~\cite{li2024chainideasrevolutionizingresearch} structures relevant literature into a chain scaffold that reflects a field’s progression and can be revisited as new evidence arises. The AI Scientist-v2~\cite{yamada2025ai} incorporates a progressive agentic tree-search approach that enables branching, backtracking and follow-up experimentation across iterations.




\paragraph{Agentic Feedback and Reflection.}


Modern web agents treat interaction as a continual learning process, using feedback signals and reflection modules to refine their reasoning and recover from failures over time. Agent Q~\cite{putta2024agent} combines guided Monte Carlo tree search with a self-critique stage, so that rollouts provide not only action sequences but also preference-style supervision. \textsc{ReAP}~\cite{azam2025reflection} makes reflection explicit by treating it as a retrieval problem: it stores task–reflection key–value pairs summarizing what was learned from past trajectories, then, at inference time, retrieves the most relevant reflections and appends them to the agent’s prompt to guide planning on new web-navigation tasks. Agent-E~\cite{azam2024multimodal} introduces an automatic validation pipeline that detects execution errors across text and vision, and then triggers self-refinement, enabling agents to iteratively correct their own workflows. Recon-Act~\cite{he2025recon} uses a dual-team architecture in which a Reconnaissance team extracts generalized tools from successful and failed trajectories, and an Action team applies these tools to re-plan tasks, forming a closed feedback loop. INFOGENT~\cite{reddy2025infogent}, on the other hand, leverages aggregator feedback to iteratively refine navigation and search strategies based on identified information gaps. And WINELL~\cite{reddy2025winell}, as a updating web agent, relies on feedback from the aggregation process to adapt subsequent searches and update selection during continuous operation. Finally, self-reflective search agents such as WebSeer~\cite{he2025webseer} integrate explicit self-reflection signals into reinforcement learning, constructing reflection-annotated trajectories and a two-stage training framework so that mis-solved or uncertain cases become targeted feedback that deepens future search and reasoning.

GUI agents also integrate explicit reflection so they can critique and repair their own plans. Early computer-control systems with structured reflection, for example, a zero-shot desktop control agent with structured self-reflection loops~\cite{li2023zero}, provides conceptual templates that later GUI agents adapt to visual, multi-application settings. GUI-Reflection~\cite{wu2025gui} instantiates this idea end-to-end: it builds a reflection-oriented task suite, automatically synthesizes error scenarios from existing successful trajectories, and adds an online reflection-tuning stage so multi-modal GUI models learn to detect failures, reason about causes, and generate corrective actions without human annotation. History-Aware Reasoning (HAR)~\cite{wang2025history} treats long-horizon GUI automation as a reflective learning problem, constructing reflective learning scenarios, synthesizing tailored correction guidelines, and designing a hybrid RL reward so the agent acquires episodic reasoning knowledge from its own errors and shifts from history-agnostic to history-aware reasoning. MobileUse~\cite{li2025mobileuse} introduces hierarchical reflection on mobile devices, where the agent self-monitors at the action, subtask, and task level and triggers reflection on demand, pausing only when needed to diagnose and recover. InfiGUIAgent~\cite{liu2025infiguiagent} integrates hierarchical and expectation–reflection reasoning in a second training stage, enabling the agent to run expectation–reflection cycles that compare expected and actual outcomes and revise multi-step plans when they diverge. Mobile-Agent-E~\cite{wang2025mobileagente} embeds an Action Reflector and Notetaker that evaluate executed steps and write refined Tips and Shortcuts back into persistent long-term memory, forming a self-evolution loop where the agent’s behavior is progressively refined from accumulated experience. 


For autonomous research agents, learning from outcomes is essential to improve reasoning and experimental reliability over time. CycleResearcher~\cite{weng2025cycleresearcherimprovingautomatedresearch} couples a research agent with a reviewer agent that provides automated peer-review feedback, and uses an iterative preference-training loop so the research agent can refine future drafts and decisions. MLR-Copilot~\cite{li2024mlrcopilotautonomousmachinelearning} monitors execution results and human comments during experiment implementation and execution, using these signals to iteratively refine code, configurations and even upstream hypotheses. Dolphin~\cite{yuan2025dolphin} implements a closed-loop auto-research framework in which generated code is run on benchmarks and exception-guided debugging plus outcome analysis feed back into idea generation and implementation, pruning unproductive paths. At the search–reasoning interface, DeepResearcher~\cite{zheng2025deepresearcher} optimizes query, browsing, and answering policies via reinforcement learning on real-web trajectories, with outcome rewards inducing behaviors such as planning, cross-validation, and self-reflection. Agentic Deep Research~\cite{weizhi2025from} further emphasizes reward design for reasoning-driven search, arguing that principled incentives over answer quality and reasoning traces provide structured signals that improve downstream synthesis in deep-research agents.




\subsubsection{Collective multi-agent reasoning}
\label{sec:app-auto-multi-agent}


% web agent
Collective multi-agent reasoning for web agents reframes browser use as cooperation among specialized roles rather than a single monolithic policy. WebPilot~\cite{zhang2025webpilot} models web task execution as a multi-agent system with a global planning agent that decomposes tasks and local MCTS-based executors that solve subtasks, jointly steering search in complex web environments. INFOGENT~\cite{reddy2025infogent} organizes web information aggregation into a Navigator, Extractor, and Aggregator, so exploration, evidence extraction, and synthesis are handled by distinct cooperating agents with feedback from the Aggregator to guide future navigation; WINELL~\cite{reddy2025winell} leverages agentic web search to plan and execute iterative information gathering for discovering timely factual updates relevant to a target Wikipedia article.. Recon-Act~\cite{he2025recon} adopts a Reconnaissance–Action paradigm in which a Recon team analyzes successful and failed trajectories to derive generalized tools or hints, and an Action team re-plans and executes with this evolving toolset. PAE~\cite{zhou2025proposer} uses three roles, namely a task proposer, an acting web agent and a VLM-based evaluator, to autonomously generate vision-based web tasks and feed success signals back into the policy via RL. Hierarchical web agents such as Agent-E~\cite{abuelsaad2024agent} and Plan-and-Act~\cite{wang2023plan} similarly separate a high-level planner from a browser-navigation agent, enabling structured plan–execute cooperation. At a more conceptual level, Agentic Web~\cite{yang2025agentic} envisions the internet as an agentic web of interacting agents and analyzes how coordination, communication protocols and economic incentives shape such ecosystems, while Agentic Deep Research~\cite{weizhi2025from} frames information seeking as iterative feedback loops of reasoning, retrieval, and synthesis that can be instantiated by single- or multi-agent web research systems.



% gui agent
Multi-agent designs for GUI agents typically decompose “using a computer” into cooperating roles that plan, perceive, decide and execute. COLA~\cite{zhao2025cola} instantiates a scenario-aware task scheduler, a planner, a decision-agent pool, an executor, and a reviewer, so UI tasks are split into basic capability units and routed to domain-specialized agents rather than a single monolith. On mobile, Mobile-Agent-v2~\cite{wang2024mobile} adopts a tri-role pattern with planner, decision, and reflection agents for progress navigation, local action selection, and error correction, while Mobile-Agent-E~\cite{wang2025mobileagente} further builds a hierarchical stack with a Manager and four subordinate agents (i.e. Perceptor, Operator, Action Reflector, Notetaker) plus a self-evolution module that learns long-term Tips and Shortcuts from experience. Mobile-Agent-V~\cite{wang2025mobileagentvvideoguidedapproacheffortless} similarly employs a video agent, decision agent, and reflection agent to coordinate multi-modal perception and execution, and MobileExperts~\cite{zhang2024mobileexperts} dynamically forms teams of expert agents with a dual-layer planner that allocates subtasks to tool-specialized experts. SWIRL~\cite{lu2025swirl} makes this structure explicit for RL, training a Navigator that converts language and screen context into structured plans and an Interactor that grounds those plans into atomic GUI actions within a multi-agent RL workflow. PC Agent~\cite{liu2025pc} uses separate planning and grounding agents in a two-stage pipeline for desktop automation, illustrating how multi-agent decomposition can improve long-horizon PC control. 

% research agent
To facilitate autonomous research agents, multi-agent collaboration enables a single model’s linear workflow to become a coordinated research group: specialized agents operate in parallel, exchange intermediate artifacts through explicit interfaces, and provide adversarial or complementary feedback to improve both creativity and rigor. For example, AgentRxiv~\cite{schmidgall2025agentrxivcollaborativeautonomousresearch} coordinates author, reviewer, and editor agents that iteratively refine manuscripts and share evolving artifacts across virtual “labs.” ARIA~\cite{ramirezmedina2025acceleratingscientificresearchmultillm} instantiates a role-structured multi-LLM team that searches, filters, and synthesizes scientific literature into actionable experimental procedures. Earlier multi-agent designs such as CAMEL~\cite{qi2023largelanguagemodelszero} demonstrate how cooperative role-play with tool access can enhance hypothesis generation and task decomposition. In experimental sciences, Coscientist~\cite{boiko2023autonomous} integrates planning, robotic instrument control, and analysis into a multi-agent closed loop that autonomously designs and executes wet-lab experiments. Finally, TAIS~\cite{liu2024toward} defines a hierarchical team, namely project manager, data engineer and domain expert, that jointly discovers disease-predictive genes from expression data through coordinated division of labor.





